% Part:sets-functions-relations
% Chapter: size-of-sets
% Section: equinumerous-sets

\documentclass[../../../include/open-logic-section]{subfiles}

\begin{document}

\olfileid{sfr}{siz}{equ}

\olsection{ہم تعدادیت}

ساڈے کول سیٹاں دے ”سائز“ دا اک بدیہی تصور اے، جیہڑا متناہی سیٹاں
لئی چنگی طرح کم کردا اے۔ پر لامتناہی سیٹاں دا کی؟ جے اسی
\emph{کسے وی} سائز دے دو سیٹاں دے سائز دا رسمی موازنہ کرن دا طریقہ
بناؤنا چاہنے آں، تاں ایہ تعریف کرن توں شروع کرنا چنگا خیال اے کہ سیٹ
کدوں اکّو سائز دے ہوندے نیں۔ فریگے دی مثال ویکھو:
\begin{quote}
  جے اک بیرا ایہ پکا کرنا چاہوے کہ اوہنے میز اُتے پلیٹاں دے عین برابر
  چھریاں رکھیاں نیں، تاں اوہنوں دوہاں وچوں کسے نوں وی گنّن دی لوڑ
  نہیں، جے اوہ بس ہر پلیٹ دے سجے پاسے اک چھری رکھ دیوے، ایس طرح کہ
  میز اُتے ہر چھری کسے پلیٹ دے سجے پاسے پئی ہووے۔ ایس طرح پلیٹاں تے
  چھریاں اک دوجے نال اکلوتے ڈھنگ نال میل کھاندیاں نیں، تے اوہ وی
  اوسے مکانی تعلق راہیں۔ \citep[\S70]{Frege1884}
\end{quote}
ایس عبارت دی سوجھ نوں اک رسمی تعریف راہیں صاف کیتا جا سکدا اے:

\begin{defn}\ollabel{comparisondef}
  $A$،~$B$ دے نال \emph{ہم تعداد} اے، جیہنوں $\cardeq{A}{B}$ لکھدے
  آں، عین اوس ویلے جدوں کوئی !!a{bijection} $f \colon A \to B$ موجود ہووے۔
\end{defn}

\begin{prop}\ollabel{equinumerosityisequi}
ہم تعدادیت اک ہم ارزیت دا تعلق اے۔
\end{prop}

\begin{proof}
سانوں وکھاؤنا اے کہ ہم تعدادیت انعکاسی، متناظر تے متعدی اے۔
$A, B$ تے~$C$ سیٹ ہون۔

\emph{انعکاسیت۔} شناختی نقشہ $\Id{A} \colon A \to A$، جتھے
$\Id{A} (x) = x$ ہر $x \in A$ لئی، !!a{bijection} اے۔ سو
$\cardeq{A}{A}$۔

\emph{تناظر۔} فرض کرو $\cardeq{A}{B}$، یعنی کوئی
!!a{bijection} $f\colon A \to B$ موجود اے۔ کیوں جے $f$،
!!{bijective} اے، ایس لئی اوہدا اُلٹ $f^{-1}$ موجود اے تے اوہ وی
!!{bijective} اے۔ ایس لئی $f^{-1}\colon B \to A$ اک
!!a{bijection} اے، سو $\cardeq{B}{A}$۔

\emph{تعدیت۔} فرض کرو کہ $\cardeq{A}{B}$ تے $\cardeq{B}{C}$، یعنی
!!{bijection}s $f\colon A \to B$ تے $g\colon B \to C$ موجود نیں۔
فیر ترکیب $\comp{f}{g}\colon A \to C$، !!{bijective} اے، سو
$\cardeq{A}{C}$۔
\end{proof}

\begin{prop}
جے $\cardeq{A}{B}$، تاں $A$، !!{enumerable} عین اوس ویلے اے جدوں
$B$ اے۔
\end{prop}

\begin{editorial}
اگلا ثبوت، جے \olref[enm]{sec} شامل ہووے، تاں
\olref[enm]{defn:enumerable} ورتدا اے، تے ہور صورت وچ
\olref[enm-alt]{defn:enumerable} ورتدا اے۔
\end{editorial}

\begin{proof}
فرض کرو $\cardeq{A}{B}$، سو کوئی !!{bijection} $f \colon A
\to B$ موجود اے، تے فرض کرو کہ $A$، !!{enumerable} اے۔
\oliflabeldef{sfr:siz:enm:defn:enumerable}{
  فیر یا تاں $A = \emptyset$ اے یا کوئی !!a{surjective} فنکشن
  $g\colon \PosInt \to A$ موجود اے۔ جے $A = \emptyset$، تاں
  $B = \emptyset$ وی اے (نہیں تاں کوئی !!a{element}~$y \in B$
  ہوندا، پر کوئی $x \in A$ نہ ہوندا جیہدے لئی $f(x) = y$)۔ جے، دوجے
  پاسے، $g\colon \PosInt \to A$، !!{surjective} اے، تاں
  $\comp{g}{f} \colon \PosInt \to B$، !!{surjective} اے۔ ایہ ویکھن
  لئی $y \in B$ لو۔ کیوں جے $f$، !!{surjective} اے، ایس لئی کوئی
  $x \in A$ اے جیہدے لئی $f(x) = y$۔ کیوں جے $g$، !!{surjective}
  اے، ایس لئی کوئی $n \in \PosInt$ اے جیہدے لئی $g(n) = x$۔ ایس لئی،
\[
(\comp{g}{f})(n) = f(g(n)) = f(x) = y
\]
تے سو $\comp{g}{f}$، !!{surjective} اے۔ ساڈے کول
$\comp{g}{f}$،~$B$ دی اک گنتی وار فہرست اے، تے ایس لئی $B$،
!!{enumerable} اے۔}
{
  فیر یا تاں $A = \emptyset$ اے یا کوئی !!a{bijection}~$g$ موجود
  اے جیہدی حاصل قدراں دا سیٹ~$A$ اے تے جیہدا دائرۂ تعریف یا $\Nat$
  اے یا قدرتی عدداں دا اک مُڈھلا سلسلہ اے۔ جے $A = \emptyset$، تاں
  $B = \emptyset$ وی اے (نہیں تاں کوئی~$y \in B$ ہوندا جیہدے لئی
  کوئی $x \in A$ ای نہ ہوندا کہ $f(x) = y$)۔ سو فرض کرو ساڈے کول
  ساڈی !!{bijection}~$g$ اے۔ فیر $\comp{g}{f}$ اک !!a{bijection} اے
  جیہدی حاصل قدراں دا سیٹ~$B$ اے تے جیہدا دائرۂ تعریف~$g$ والے دے
  برابر اے (یعنی یا $\Nat$ یا اوہدا اک مُڈھلا ٹکڑا)، سو $B$،
  !!{enumerable} اے۔}
% OLSIZ-006 ماخذ دی درستی: خالی سیٹ والیاں دونیں مشروط شاخاں وچ f ہی زیرِ بحث مطابقت اے؛ انگریزی ماخذ اوتھے g(x)=y اوہدے تعارف توں پہلاں لکھدا سی۔

جے $B$، !!{enumerable} اے، تاں اسی حاصل کردے آں کہ $A$،
!!{enumerable} اے، دلیل نوں !!{bijection} $f^{-1}\colon B \to A$
نال دُہرا کے،~$f$ دی بجائے۔
\end{proof}

\begin{prob}
وکھاؤ کہ جے $\cardeq{A}{C}$ تے $\cardeq{B}{D}$، تے $A \cap B =
C \cap D = \emptyset$، تاں $\cardeq{A \cup B}{C \cup D}$۔
\end{prob}

\begin{prob}
وکھاؤ کہ جے $A$ لامتناہی تے !!{enumerable} اے، تاں
$\cardeq{A}{\Nat}$۔
\end{prob}

\end{document}
